\section{Introduction}
Robotic manipulation \cite{mason2018toward} is a fundamental task in robotics, enabling robots to physically interact with their environment across a wide range of applications, including manufacturing, logistics, and autonomous systems. Manipulation strategies can be broadly classified into two categories: prehensile and non-prehensile manipulation, depending on the nature of the contact interactions between the robot and the manipulated object. While prehensile manipulation \cite{bicchi2000robotic} involves establishing and maintaining a stable grasp on an object, typically through multiple contact points, non-prehensile manipulation \cite{lynch1996nonprehensile,lynch1999dynamic} achieves object motion only through contact interactions such as pushing, sliding, and rolling. Among these techniques, pushing \cite{mason1986mechanics} is widely employed as a simple yet effective manipulation primitive for achieving desired object poses in planar environments. Unlike grasp-based approaches, pushing does not require form or force closure, making it particularly suitable for manipulating objects that are challenging to grasp. 

While single-robot pushing provides a simple and effective means of object manipulation, its capabilities may be constrained by the limited force capacity and available contact locations of an individual robot. To address these limitations, multi-robot pushing \cite{ozaki1997distributed,tuci2018cooperative} employs a team of robots to cooperatively manipulate an object by applying forces at different contact locations, overcoming the limitations of individual robots in terms of force capacity and manipulation flexibility. Rosenfelder et al. \cite{rosenfelder2024force} developed a unified framework for cooperative pushing, integrating formation optimization, force allocation, and distributed control. This comprehensive approach provides a systematic solution for coordinating multiple robots to achieve reliable and efficient object transportation.

The inertial parameters of an object, particularly its CoM, play a crucial role in determining its dynamic response to applied pushing forces. While some adaptive pushing control methods do not require prior knowledge of the object's CoM \cite{wang2016kinematic,nguyen2026dnpush}, accurate estimation of its location can provide valuable insight into the object's physical characteristics, facilitating more informed and reliable manipulation planning. In particular, CoM information can guide the determination of appropriate contact configurations and pushing strategies to produce the desired object motion. For instance, Cosgun et al. \cite{cosgun2011push} employed the object's CoM as a reference for initializing gripper placement when planning linear pushing actions. Similarly, Nguyen et al. \cite{nguyen2024distributed} introduced a distributed cascade force control strategy for cooperative object transportation, assuming that the robots are positioned on opposite sides of the projection of the object's CoM onto the contact edge. 

Despite its importance in manipulation planning and control, the object's CoM is often assumed to be known or approximated from geometric information under a uniform mass distribution. However, these simplifications may not adequately capture the object's true characteristics, potentially compromising the accuracy of CoM estimates.
Conventional estimation methods typically rely on dynamic models derived from the Newton–Euler equations \cite{Khosla1985Parameter,Swevers2007Dynamic,Gautier1997Dynamic,Dong2018Payload,Mellinger2011Aerial,Lee2017AerialTransportation,Yu2005Pushing,Franchi2014Distributed,Franchi2016Decentralized}, which relate applied forces and torques to the resulting linear and angular accelerations, which are susceptible to measurement noise and uncertainties that may degrade the accuracy and robustness of the estimation process. Recent studies have investigated alternative control-based approaches that enable CoM estimation through deliberately induced object motions, circumventing the need for explicit dynamic modeling \cite{Nowak2008Generalized,Hyland2023CoM,Hyland2025Onboard,Uemura2024Dynamic}. However, these approaches often require multiple pushing actions and are primarily designed for single-robot manipulation, limiting their applicability to multi-robot systems. 
To address this problem, we propose a distributed strategy that enables a team of omnidirectional mobile robots to estimate the CoM of a planar polygonal object using only pushing contacts. The key idea is to coordinate the robots such that the net pushing force on the object vanishes while a nonzero torque is generated, inducing an approximately pure rotation about the CoM. The main contributions of this paper are as follows:
\begin{itemize}
    \item A formation synthesis and force control framework that generates a zero-net-force, rotation-inducing excitation for the object, with the formation optimized for robustness using ALPSO.
    \item A distributed RLS-based estimator that allows each robot to estimate the CoM independently from local kinematic measurements only, without requiring acceleration measurements or information exchange.
    \item Numerical simulations demonstrating millimeter-level estimation accuracy within a single pushing operation under measurement noise.
\end{itemize}

The remainder of this paper is organized as follows. Section~\ref{sec:related} reviews related work on inertial parameter and CoM estimation. Section~\ref{sec:method} formulates the problem and presents the estimation algorithm, formation synthesis, and force control. Section~\ref{sec:results} reports the simulation results, and Section~\ref{sec:conclusion} concludes the paper.